\chapter{Authentication and authorization}

Authn proves who a user is.
Authz decides what they are allowed to do.
Mixing them creates security gaps.

\section{Authentication}
Common strategies include sessions, API keys, and JWTs.
Pick one and document its lifecycle and rotation policy.

\section{Authorization}
Authorization should be explicit and testable.
Use roles or permissions and keep them small.

\section{Resource ownership}
Many systems can model access as ownership or membership.
Capture it in the data model and enforce it at the boundary.

\section{Service-to-service auth}
Use short-lived tokens and mutual TLS when possible.
Avoid long-lived static keys.

\begin{warningbox}{Warning: auth in the UI is not enough}
Never rely on frontend checks to enforce permissions.
Always enforce authorization on the server.
\end{warningbox}

\section{Token lifecycle}
Tokens should be short-lived and rotated regularly.
If you use refresh tokens, store them securely and revoke them on logout.
Document token expiration and refresh behavior for clients.

\section{Session management}
Session-based auth requires careful cookie settings.
Use \code{HttpOnly}, \code{Secure}, and appropriate \code{SameSite} values.
Invalidate sessions on password resets or suspicious activity.

\section{Claims and scopes}
JWT claims should be minimal and stable.
Scopes are easier to reason about than complex roles.
Keep a mapping table between scopes and allowed actions.

\section{RBAC and ABAC}
Role-based access control (RBAC) is easy to operate.
Attribute-based access control (ABAC) is more flexible but harder to test.
Choose the simplest model that meets your needs.

\section{Resource checks}
Authorization should happen before expensive work.
Check ownership or membership before you query large datasets.
This reduces data leakage and improves performance.

\section{Service identity}
Treat service identities like users.
Each service should have its own identity and permissions.
Do not share credentials across services.

\section{Audit trails}
Record authentication and authorization events.
Audit logs help you answer who accessed what and when.
Keep audit logs immutable and protected.

\section{Key rotation}
Keys and secrets should rotate regularly.
Build rotation into the system so it is routine, not an emergency.
If rotation requires downtime, the design is fragile.

\section{Example policy table}
\begin{tabularx}{\textwidth}{@{}lX@{}}
\toprule
Action & Required scope \\
\midrule
Read tasks & \code{tasks:read} \\
Create tasks & \code{tasks:write} \\
Delete tasks & \code{tasks:admin} \\
\bottomrule
\end{tabularx}

\section{Error responses}
Return \code{401} for missing or invalid credentials.
Return \code{403} for valid credentials without permission.
Do not leak resource existence in authorization errors.

\section{Password storage}
If you store passwords, never store them in plain text.
Use modern hashing algorithms and a strong cost factor.
Rotate hashing parameters when industry guidance changes.

\section{MFA and step-up auth}
Some actions need stronger authentication.
Use step-up authentication for sensitive operations like deletions and billing changes.
This reduces the impact of stolen tokens.

\section{Delegated access}
If users authorize third-party access, use explicit scopes and expiry.
OAuth-style flows are common for this scenario.
Always allow users to revoke access.

\section{API keys}
API keys should be scoped and revocable.
Avoid user-level keys for server-to-server traffic.
Log key usage so you can detect abuse.

\section{Rate limiting}
Auth endpoints are high-value targets.
Apply stricter rate limits to login and token refresh endpoints.
Return clear errors and do not leak which part failed.

\section{Policy evaluation}
Authorization checks should be centralized.
If each endpoint rolls its own policy, rules will drift.
Consider a policy engine when rules become complex.

\section{Admin actions}
Admin operations need extra logging and review.
Require elevated scopes and generate audit events.
If admin actions are not tracked, incidents are inevitable.

\section{Service auth example}
\begin{lstlisting}[language=JSON]
{
  "service_id": "billing-worker",
  "scopes": ["tasks:read", "tasks:write"],
  "issued_at": "2026-02-18T10:00:00Z"
}
\end{lstlisting}

\section{Secret storage}
Store secrets in a dedicated secret manager.
Do not store secrets in code repositories or environment files.
Access should be limited to the services that need them.

\section{Threat model}
List the threats you care about:
\begin{itemize}[nosep]
  \item credential stuffing
  \item token replay
  \item privilege escalation
  \item insider access
\end{itemize}

The threat model guides which controls you implement first.

\section{Identity proofing}
Not all users are equal risk.
For high-risk flows, add identity proofing such as verified email or phone.
For enterprise accounts, consider domain verification before granting admin access.

\section{Account lifecycle}
Accounts have states: active, locked, disabled, deleted.
Define how each state affects authentication and authorization decisions.
State transitions should be auditable and reversible where possible.

\section{Password resets}
Reset tokens should be short-lived and single use.
Do not reveal whether an email exists in the system.
Send reset alerts and require re-authentication after password changes.

\section{Session binding}
Bind sessions to device or IP when risk warrants it.
If the context changes significantly, require re-authentication.
This mitigates stolen session cookies.

\section{Token revocation}
Short-lived tokens reduce risk, but revocation still matters.
Maintain a revocation list for compromised tokens.
If you cannot revoke, shorten expiry and rely on rotation.

\section{Scope design}
Keep scopes small and action-oriented.
Avoid scopes that combine read and write unless necessary.
Document scope usage with examples so clients can request the minimal set.

\section{Fine-grained checks}
Role checks alone are not sufficient for multi-tenant systems.
Combine role checks with resource ownership checks.
This avoids privilege escalation within the same role.

\section{Policy composition}
Complex policies should be composed, not duplicated.
Create reusable policy functions like \code{can\_manage\_project}.
Composition keeps policy logic consistent across endpoints.

\section{Authorization at the data layer}
Defense in depth is valuable.
If possible, enforce row-level security or data filters in the database.
This reduces the risk of missing checks in the application layer.

\section{User impersonation}
Support impersonation only with strict controls.
Require explicit justification and log every impersonation action.
Impersonation sessions should have short expiry.

\section{Audit log structure}
Audit logs should include actor, action, resource, and outcome.
Use immutable storage or write-once retention when required.
Include correlation ids so incidents can be traced end-to-end.

\section{Key management}
Separate keys by environment and purpose.
Use asymmetric keys for tokens when possible.
Rotate keys on a schedule and on incident response triggers.

\section{Secret sprawl}
Track where secrets are used.
Rotate and prune secrets that are no longer referenced.
Secret sprawl is a hidden security risk.

\section{Service identity onboarding}
Provision service identities using automation, not manual steps.
Automated provisioning reduces mistakes and missing permissions.
Store service identity metadata with owners and expiry dates.

\section{Security testing}
Add tests for authorization boundaries.
Use negative tests that assert access is denied.
Security tests should run in CI, not just as manual checks.

\section{Auth error responses}
Return errors that are informative without leaking sensitive details.
Avoid differences between invalid user and invalid password responses.
Keep response times consistent to reduce timing attacks.

\section{Example policy evaluation flow}
\begin{lstlisting}[language=Python]
def authorize(user, action, resource):
    if not user.is_authenticated:
        return "unauthenticated"
    if not has_scope(user, action):
        return "forbidden"
    if not owns_resource(user, resource):
        return "forbidden"
    return "ok"
\end{lstlisting}

\section{Incident response}
Document how to revoke access quickly.
Predefine emergency procedures for compromised keys or accounts.
Fast response reduces the blast radius of auth failures.

\section{Single sign-on}
Enterprise users often require SSO.
Support SAML or OIDC flows with clear documentation.
Test SSO flows in staging with real identity providers.

\section{Provisioning and deprovisioning}
User provisioning should be automated when possible.
Deprovisioning must be fast and reliable.
If a user leaves, access should be revoked across all services.

\section{Access reviews}
Periodic access reviews catch permission creep.
Review roles and scopes for privileged users.
Automate review reports to reduce manual work.

\section{Device trust}
Some systems require trusted devices.
Bind sessions to devices and require re-approval when devices change.
Device trust helps reduce account takeover risk.

\section{Consent and transparency}
When access is delegated, capture consent with timestamps.
Let users review and revoke consent in a self-service UI.
Consent history helps with compliance and user trust.

\section{Token format stability}
Changing token formats can break clients.
If you must change them, version tokens explicitly.
Provide a clear migration window for old tokens.

\section{JWT headers and key ids}
Include a key id so verifiers can select the right public key.
Rotate keys and publish a key set endpoint.
Avoid manual key distribution.

\section{mTLS for internal traffic}
Mutual TLS provides strong service identity.
It can reduce the need for API keys in internal networks.
Use automation for certificate rotation.

\section{Authorization caching}
Authorization checks can be cached for short windows.
Cache only when permissions change infrequently.
Invalidate caches on role changes to avoid stale access.

\section{Auth metrics}
Track login success rates, token refresh errors, and authorization denials.
These metrics reveal attacks and misconfigurations.
Alert on unusual spikes in failed auth attempts.

\section{Policy drift}
As systems evolve, policy logic drifts.
Keep policy tests in a dedicated suite.
If policy logic changes, update the policy documentation too.

\section{Example access review checklist}
\begin{itemize}[nosep]
  \item review admin users
  \item review service account scopes
  \item remove unused API keys
  \item confirm MFA enrollment for privileged roles
\end{itemize}

\section{Tenant isolation}
Multi-tenant systems require strict tenant isolation.
Always include tenant checks in authorization rules.
Treat missing tenant checks as security bugs.

\section{Account lockout}
Lockouts can stop credential stuffing but can also be abused.
Use progressive delays and risk scoring.
Allow lockout recovery with verified identity.

\section{MFA recovery}
Provide recovery codes for MFA.
Recovery codes should be single-use and revocable.
Encourage users to rotate recovery codes after use.

\section{Session limits}
Limit concurrent sessions for sensitive accounts.
Expose session management in user settings.
Let users revoke sessions from a trusted device.

\section{Policy as code}
Represent policies in code with tests and code review.
Policy changes should follow the same deployment process as other code.
Avoid editing policy rules directly in production.

\section{Credential stuffing defense}
Use IP reputation, rate limits, and breached password detection.
Combine signals rather than relying on a single control.
This reduces both false positives and false negatives.

\section{Example auth metrics}
\begin{lstlisting}[language=YAML]
metrics:
  login_success_rate: 0.98
  token_refresh_failures: 12
  authz_denied_rate: 0.03
\end{lstlisting}

\section{Tenant admin boundaries}
Tenant admins should be limited to their tenant.
Cross-tenant access should require a separate global role.
This prevents accidental lateral access.

\section{SCIM provisioning}
SCIM enables automated user lifecycle management.
If you support SCIM, test both create and deprovision flows.
Ensure SCIM actions trigger the same audit logs as manual changes.

\section{API access versus user access}
Service API access often needs different policies than user access.
Separate scopes for machine traffic and human traffic.
This reduces the risk of over-privileged service accounts.

\section{Session timeout policy}
Define idle and absolute timeouts for sessions.
Idle timeouts reduce exposure for unattended sessions.
Absolute timeouts prevent indefinite session lifetimes.

\section{Key rotation playbook}
Publish a rotation playbook with clear steps and owners.
Practice rotations to avoid surprises during incidents.
Rotation should be a routine operation, not an emergency.

\section{Break-glass access}
Emergency access should be rare and controlled.
Use a separate break-glass role with short-lived access.
Every break-glass action should create an audit event.

\section{Just-in-time access}
Grant elevated permissions only when needed.
Just-in-time access reduces the risk of persistent over-privilege.
Automate expiry and notify owners when access is granted.

\section{Session anomaly detection}
Detect unusual login patterns such as impossible travel.
Flag anomalies for re-authentication or step-up checks.
Anomaly detection reduces account takeover risk.

\section{Security headers}
Set security headers to protect session cookies and tokens.
Use \code{Strict-Transport-Security}, \code{Content-Security-Policy}, and \code{X-Content-Type-Options}.
Headers are part of the auth surface, not just frontend hygiene.

\section{Example conditional policy}
\begin{lstlisting}[language=YAML]
policy:
  action: "export_data"
  require:
    - scope: "data:export"
    - mfa: true
    - business_hours: true
\end{lstlisting}

\begin{checklistbox}{Checklist: auth}
\begin{itemize}[nosep]
  \item Separate authentication and authorization.
  \item Keep permissions minimal and explicit.
  \item Rotate credentials regularly.
\end{itemize}
\end{checklistbox}
